\documentclass{amsart}
% For dealing with references we use the comment environment
\usepackage{verbatim}
\newenvironment{reference}{\comment}{\endcomment}
%\newenvironment{reference}{}{}
\newenvironment{slogan}{\comment}{\endcomment}
\newenvironment{history}{\comment}{\endcomment}

% For commutative diagrams we use Xy-pic
\usepackage[all]{xy}

% We use 2cell for 2-commutative diagrams.
\xyoption{2cell}
\UseAllTwocells

% We use multicol for the list of chapters between chapters
\usepackage{multicol}

% This is generally recommended for better output
\usepackage{lmodern}
\usepackage[T1]{fontenc}

% For cross-file-references

% Package for hypertext links:
\usepackage{hyperref}

% For any local file, say "hello.tex" you want to link to please

% Theorem environments.
%
\theoremstyle{plain}
\newtheorem{theorem}[subsection]{Theorem}
\newtheorem{proposition}[subsection]{Proposition}
\newtheorem{lemma}[subsection]{Lemma}

\theoremstyle{definition}
\newtheorem{definition}[subsection]{Definition}
\newtheorem{example}[subsection]{Example}
\newtheorem{exercise}[subsection]{Exercise}
\newtheorem{situation}[subsection]{Situation}

\theoremstyle{remark}
\newtheorem{remark}[subsection]{Remark}
\newtheorem{remarks}[subsection]{Remarks}

\numberwithin{equation}{subsection}

% Macros
%
\def\lim{\mathop{\mathrm{lim}}\nolimits}
\def\colim{\mathop{\mathrm{colim}}\nolimits}
\def\Spec{\mathop{\mathrm{Spec}}}
\def\Hom{\mathop{\mathrm{Hom}}\nolimits}
\def\Ext{\mathop{\mathrm{Ext}}\nolimits}
\def\SheafHom{\mathop{\mathcal{H}\!\mathit{om}}\nolimits}
\def\SheafExt{\mathop{\mathcal{E}\!\mathit{xt}}\nolimits}
\def\Sch{\mathit{Sch}}
\def\Mor{\mathop{\mathrm{Mor}}\nolimits}
\def\Ob{\mathop{\mathrm{Ob}}\nolimits}
\def\Sh{\mathop{\mathit{Sh}}\nolimits}
\def\NL{\mathop{N\!L}\nolimits}
\def\CH{\mathop{\mathrm{CH}}\nolimits}
\def\proetale{{pro\text{-}\acute{e}tale}}
\def\etale{{\acute{e}tale}}
\def\QCoh{\mathit{QCoh}}
\def\Ker{\mathop{\mathrm{Ker}}}
\def\Im{\mathop{\mathrm{Im}}}
\def\Coker{\mathop{\mathrm{Coker}}}
\def\Coim{\mathop{\mathrm{Coim}}}

% Boxtimes
%
\DeclareMathSymbol{\boxtimes}{\mathbin}{AMSa}{"02}

%
% Macros for moduli stacks/spaces
%
\def\QCohstack{\mathcal{QC}\!\mathit{oh}}
\def\Cohstack{\mathcal{C}\!\mathit{oh}}
\def\Spacesstack{\mathcal{S}\!\mathit{paces}}
\def\Quotfunctor{\mathrm{Quot}}
\def\Hilbfunctor{\mathrm{Hilb}}
\def\Curvesstack{\mathcal{C}\!\mathit{urves}}
\def\Polarizedstack{\mathcal{P}\!\mathit{olarized}}
\def\Complexesstack{\mathcal{C}\!\mathit{omplexes}}
% \Pic is the operator that assigns to X its picard group, usage \Pic(X)
% \Picardstack_{X/B} denotes the Picard stack of X over B
% \Picardfunctor_{X/B} denotes the Picard functor of X over B
\def\Pic{\mathop{\mathrm{Pic}}\nolimits}
\def\Picardstack{\mathcal{P}\!\mathit{ic}}
\def\Picardfunctor{\mathrm{Pic}}
\def\Deformationcategory{\mathcal{D}\!\mathit{ef}}

\setlength{\emergencystretch}{3em}
\begin{document}
\title[Stacks Project, Tag 05LT]{362 --- Finite-stage descent of flatness along a filtered colimit}
\maketitle
\noindent Copyright (C) 2005--2025 Johan de Jong. Stacks Project authors. Source: \href{https://stacks.math.columbia.edu/tag/05LT}{Tag 05LT}.

\noindent Editorial difficulty note (not author text). Difficulty H2: The proof extracts finitely many principal opens covering the relevant closed subset of the limit spectrum. It descends both a unit-ideal equation and local flatness to one common finite stage, then extends the argument from finite presentation to essential finite presentation..

\noindent Editorial provenance note (not author text). History: excerpt from revision \texttt{a0444\allowbreak 6e57e\allowbreak c1fbc\allowbreak 252a8\allowbreak 71afc\allowbreak ec775\allowbreak 2fb28\allowbreak 07b14}, more-algebra.tex, lines 4294--4367. Reference presentation was adapted; original-author context and core support blocks were retained or added. The mathematical wording is unchanged. Licensed under GNU FDL 1.2 or later; no invariant sections or cover texts. See ../licenses/COPYING. Context blocks reproduce author definitions and supporting results; they are not additional counted problems.

\begin{situation}
\label{situation-flattening-general}
Let $R \to S$ be a ring map. Let $J \subset S$ be an ideal.
Let $M$ be an $S$-module.
\end{situation}

In this situation, given an $R$-algebra $R'$ and an ideal $I' \subset R'$
we set $S' = S \otimes_R R'$ and $M' = M \otimes_R R'$.
We will consider the condition
\begin{equation}
\label{equation-flat-at-primes}
\forall \mathfrak q' \in V(I'S' + JS') \subset \Spec(S') :
M'_{\mathfrak q'}\text{ is flat over }R'.
\end{equation}
Geometrically, this means that $M'$ is flat over $R'$ along the intersection
of the inverse image of $V(I')$ with the inverse image of $V(J)$.
Since $(R \to S, J, M)$ are fixed, condition (\ref{equation-flat-at-primes})
only depends on the pair $(R', I')$ where $R'$ is viewed as an $R$-algebra.


\begin{lemma}
\label{lemma-limit-preserving-flat-at-primes}
In Situation \ref{situation-flattening-general}
assume $R \to S$ is essentially of finite presentation
and $M$ is an $S$-module of finite presentation. Let
$R' = \colim_{\lambda \in \Lambda} R_\lambda$
be a directed colimit of $R$-algebras. Let $I_\lambda \subset R_\lambda$
be ideals such that $I_\lambda R_\mu \subset I_\mu$ for all
$\mu \geq \lambda$ and set $I' = \colim_\lambda I_\lambda$.
If (\ref{equation-flat-at-primes}) holds for
$(R', I')$, then there exists a $\lambda \in \Lambda$ such that
(\ref{equation-flat-at-primes}) holds for $(R_\lambda, I_\lambda)$.
\end{lemma}

\begin{proof}
We first prove the lemma in case $R \to S$ is of finite presentation
and then we explain what needs to be changed in the general case.
We are going to write $S_\lambda = S \otimes_R R_\lambda$,
$S' = S \otimes_R R'$, $M_\lambda = M \otimes_R R_\lambda$, and
$M' = M \otimes_R R'$.
The base change $S'$ is of finite presentation over $R'$ and
$M'$ is of finite presentation over $S'$ and similarly for the
versions with subscript $\lambda$, see
Algebra, Lemma \href{https://stacks.math.columbia.edu/tag/05G5}{lemma base change finiteness (Tag 05G5)}.
By
Algebra, Theorem \href{https://stacks.math.columbia.edu/tag/00RC}{theorem openness flatness (Tag 00RC)}
the set
$$
U' = \{\mathfrak q' \in \Spec(S') \mid
M'_{\mathfrak q'}\text{ is flat over }R'\}
$$
is open in $\Spec(S')$. Note that $V(I'S' + JS')$
is a quasi-compact space which is contained in $U'$ by assumption.
Hence there exist finitely many $g'_j \in S'$, $j = 1, \ldots, m$
such that $D(g'_j) \subset U'$ and such
that $V(I'S' + JS') \subset \bigcup D(g'_j)$.
Note that in particular $(M')_{g'_j}$ is a flat module over $R'$.

\medskip\noindent
We are going to pick increasingly large elements $\lambda \in \Lambda$.
First we pick it large enough so that we can find
$g_{j, \lambda} \in S_{\lambda}$ mapping to $g'_j$.
The inclusion $V(I'S' + JS') \subset \bigcup D(g'_j)$ means that
$I'S' + JS' + (g'_1, \ldots, g'_m) = S'$ which can be expressed as
$$
1 = \sum y_tk_t + \sum z_sh_s + \sum f_jg'_j
$$
for some $z_s \in I'$, $y_t \in J$, $k_t, h_s, f_j \in S'$.
After increasing $\lambda$ we may assume such an equation holds in
$S_\lambda$. Hence we may assume that
$V(I_\lambda S_\lambda + J S_\lambda) \subset \bigcup D(g_{j, \lambda})$. By
Algebra, Lemma \href{https://stacks.math.columbia.edu/tag/02JO}{lemma flat finite presentation limit flat (Tag 02JO)}
we see that for some sufficiently large $\lambda$ the modules
$(M_\lambda)_{g_{j, \lambda}}$ are flat over $R_\lambda$.
In particular the module $M_\lambda$ is flat over $R_\lambda$
at all the primes corresponding to points of
$V(I_\lambda S_\lambda + J S_\lambda)$.

\medskip\noindent
In the case that $S$ is essentially of finite presentation, we can write
$S = \Sigma^{-1}C$ where $R \to C$ is of finite presentation and
$\Sigma \subset C$ is a multiplicative subset. We can also write
$M = \Sigma^{-1}N$ for some finitely presented $C$-module $N$, see
Algebra, Lemma \href{https://stacks.math.columbia.edu/tag/05N5}{lemma construct fp module (Tag 05N5)}.
At this point we introduce $C_\lambda$, $C'$, $N_\lambda$, $N'$. Then in
the discussion above we obtain an open $U' \subset \Spec(C')$
over which $N'$ is flat over $R'$. The assumption that
(\ref{equation-flat-at-primes}) is true means that $V(I'S' + JS')$ maps
into $U'$, because for a prime $\mathfrak q' \subset S'$, corresponding
to a prime $\mathfrak r' \subset C'$ we have
$M'_{\mathfrak q'} = N'_{\mathfrak r'}$. Thus we can find
$g'_j \in C'$ such that $\bigcup D(g'_j)$ contains the image of
$V(I'S' + JS')$. The rest of the proof is exactly the same as before.
\end{proof}
\end{document}
